\begin{afterword}
\summaryhead

The post continues from ``Universal Fire.'' Lavoisier united breathing and burning; before Lavoisier, Newton united the planets and the falling apple, and from Newton's work came the idea of a universal law, the same everywhere and always, with no exceptions.

To human minds this idea seems fanatical. Everyday categories are leaky, and even moral rules, the one place we want rules to hold for everyone, give way under threat. But the universe is, by human standards, ``insanely rigid.'' As far as we know, conservation of momentum has never been violated. When a scientific model fails, its replacement is again exceptionless, and the universe was obeying it all along, as general relativity governed Mercury's orbit before anyone knew it. The author gives 80 per cent confidence that the lightspeed limit is absolutely inviolable, which is not the same as thinking it holds 80 per cent of the time. The Greeks, who thought contrived experiments gave ``monstrous'' results, missed this. The post ends with three lines: since the beginning, not one unusual thing has ever happened.

\responsehead

The post's core is sound. The difference between a model and the law it describes, and between uncertainty about whether a law is exceptionless and a law with exceptions, is correct. The post also marks the limit of its own confidence: that the law was always constant is what ``our new model tells us.''

The history is compressed. The idea of a universal law did not start with Newton's discovery. Descartes, writing between 1629 and 1633, called the rules by which matter changes ``the `laws of nature','' and Newton's own belief in absolute, universal laws was shaped by a religious tradition older than both. Newton's gravitation gave the idea its most striking example; it did not create it. The unification credited to Lavoisier was also older, as the notes on ``Universal Fire'' show.

The claim about the Greeks is given with no source. Some Greek work sits uneasily with it. Galen, in the second century, tied off nerves and ureters in living animals to find out how voice and kidneys work, and took the results as a guide to the human body.

The closing lines are offered as ``the Tao'' with no attribution. I could not find them before this post, so they appear to be the author's own.

In short: a correct distinction between laws, models and our confidence in them, set in a history that dates the idea of universal law to Newton, and a claim about Greek science given without a source.
\end{afterword}
